
% Montaje experimental cinemático.
% Sin entorno figure, caption ni label.

\begin{tikzpicture}[
    >=Stealth,
    line cap=round,
    line join=round,
    font=\small,
    x=1cm,
    y=1cm
]

% =================================================
% 1. RIEL
% =================================================

\draw[line width=1.2pt] (-0.8,0) -- (10.8,0);

\node[anchor=north] at (5,-0.17)
    {Riel de $2\,\mathrm{m}$};

% =================================================
% 2. BLOQUES DE LOS EXTREMOS
% =================================================

\draw[
    fill=black!16,
    draw=black,
    line width=0.8pt
] (-0.15,0) rectangle (0.15,0.72);

\draw[
    fill=black!16,
    draw=black,
    line width=0.8pt
] (9.85,0) rectangle (10.15,0.72);

% =================================================
% 3. MICROFONOS SOBRE LOS BLOQUES
% =================================================

% Micrófono 1
\draw[line width=1pt,black!65]
    (0,0.72) -- (0,1.62);

\node[
    circle,
    draw=black!45,
    fill=black!12,
    line width=0.9pt,
    minimum size=6mm,
    inner sep=0pt
] at (0,1.62) {};

\node[anchor=south] at (0,1.94)
    {Mic. 1};

% Micrófono 2
\draw[line width=1pt,black!65]
    (10,0.72) -- (10,1.62);

\node[
    circle,
    draw=black!45,
    fill=black!12,
    line width=0.9pt,
    minimum size=6mm,
    inner sep=0pt
] at (10,1.62) {};

\node[anchor=south] at (10,1.94)
    {Mic. 2};

% =================================================
% 4. CAUCHOS / RESORTES DE RETORNO
% =================================================

% Resorte izquierdo
\draw[
    red!75!black,
    line width=1.5pt
]
    (0.15,0.48)
    -- (0.22,0.48)
    -- (0.27,0.57)
    -- (0.33,0.39)
    -- (0.39,0.57)
    -- (0.45,0.39)
    -- (0.51,0.57)
    -- (0.57,0.39)
    -- (0.63,0.48)
    -- (0.70,0.48);

% Resorte derecho
\draw[
    red!75!black,
    line width=1.5pt
]
    (9.30,0.48)
    -- (9.37,0.48)
    -- (9.43,0.57)
    -- (9.49,0.39)
    -- (9.55,0.57)
    -- (9.61,0.39)
    -- (9.67,0.57)
    -- (9.73,0.39)
    -- (9.79,0.48)
    -- (9.85,0.48);

% =================================================
% 5. CARRITO MOVIL
% =================================================

% Cuerpo bajado 0.11 cm
\draw[
    rounded corners=2pt,
    fill=black!4,
    draw=black,
    line width=0.85pt
] (1.35,0.27) rectangle (2.65,0.87);

% Ruedas: ahora su punto inferior coincide con y=0
\draw[
    fill=white,
    draw=black,
    line width=0.7pt
] (1.65,0.105) circle (0.105);

\draw[
    fill=white,
    draw=black,
    line width=0.7pt
] (2.35,0.105) circle (0.105);

\node at (2.00,0.57)
    {Carrito};

% =================================================
% 6. BUZZER EN EL EXTREMO IZQUIERDO DEL CARRITO
% =================================================

% El buzzer conserva la altura de los micrófonos.
% Se ajusta la longitud del soporte para conectarlo
% con el cuerpo del carrito.

\draw[line width=0.85pt]
    (1.55,0.87) -- (1.55,1.47);

\node[
    circle,
    draw=orange!85!black,
    fill=orange!18,
    line width=1pt,
    minimum size=6mm,
    inner sep=0pt
] at (1.55,1.62) {};

\node[anchor=west] at (1.82,1.62)
    {Buzzer};

% =================================================
% 7. LIBERACION Y RECORRIDO
% =================================================

\node[anchor=south] at (1.55,2.32)
    {Liberación};

\draw[
    ->,
    line width=0.85pt
]
    (1.55,2.27) -- (1.55,1.96);

\draw[
    ->,
    line width=1pt
]
    (2.85,2.32) -- (8.15,2.32);

\node[
    fill=white,
    inner sep=3pt
]
    at (5.50,2.32)
    {Recorrido libre};

% =================================================
% 8. DISTANCIA ENTRE MICROFONOS: 200 cm
% =================================================

\draw[densely dashed,black!45]
    (0,-0.05) -- (0,-0.85);

\draw[densely dashed,black!45]
    (10,-0.05) -- (10,-0.85);

\node[
    fill=white,
    inner sep=2pt
]
    at (0,-0.58)
    {$0\,\mathrm{cm}$};

\node[
    fill=white,
    inner sep=2pt
]
    at (10,-0.58)
    {$200\,\mathrm{cm}$};

\draw[
    <->,
    line width=0.85pt
]
    (0,-0.85) -- (10,-0.85);

\node[
    fill=white,
    inner sep=3pt
]
    at (5,-1.10)
    {Separación entre micrófonos: $200\,\mathrm{cm}$};

% =================================================
% 9. LONGITUD UTIL: 180 cm
% =================================================

\draw[densely dashed,black!45]
    (0.70,-0.53) -- (0.70,-1.55);

\draw[densely dashed,black!45]
    (9.30,-0.53) -- (9.30,-1.55);

\draw[
    <->,
    line width=0.85pt
]
    (0.70,-1.43) -- (9.30,-1.43);

\node[
    fill=white,
    inner sep=3pt
]
    at (5,-1.70)
    {Longitud útil aproximada: $180\,\mathrm{cm}$};

% =================================================
% 10. PROTOCOLO EXPERIMENTAL
% =================================================

\node[
    draw=black!65,
    rounded corners=2pt,
    fill=black!2,
    line width=0.65pt,
    text width=3.65cm,
    align=left,
    inner sep=9pt,
    anchor=west
] at (11.05,1.45) {%
    \textbf{Protocolo previsto}\\[6pt]
    10 ensayos\\[3pt]
    Masa nominal: $200\,\mathrm{g}$\\[3pt]
    Adquisición: $10\,\mathrm{s}$
};

\end{tikzpicture}
